\section{The lower determinant estimate}\label{c3:section}

Let $W$ be the good subspace on $Y_0$ introduced in
\S\ref{sz:setup}. The archimedean factor at $Y_0$
is the prescribed discrete-series type. The algebraic Weyl-valued
model used to measure a lifted vector is on the definite space $Y_m$.
We use the common $H,U,V,L$ and write
\[
 B_0=U/(mk),\qquad
 A_0=\sum_{\tau\in\Phi}
   \bigl(j(\kappa_\tau)-
          S(\kappa_\tau,\tau)\bigr),\qquad H=mk A_0.
\]
With $j$ and $S$ as in \eqref{sz:height},
$k=r+1-m$, and the difference of ranks in a single lift is
$\delta=5r$. The lift of the ordered family will always be the full
chain $\Theta_\eta$. For an $m$-fold twist $\eta$, set
$N_\eta=\prod_{w\in Q(\eta)}q_w$, so $N_1=1$. At a prime $\ell$
the allowed twists satisfy $(\ell,N_\eta)=1$.
Theorem C, below, uses the construction hypotheses in
\S\ref{c3:frontier}.

\paragraph{Local and archimedean structures.}
Use the identifications, integral modules, coordinate vector $v_1(w)$ of $\Theta_1w$, and rational metric $G_{\rm cl}$ of \S\ref{sz:setup}. Write
$\bar\iota=\iota\circ\mathrm{conj}$. The chosen
level makes the stabilizers of the final class representatives trivial.
Indeed, if a rational unitary element $u$ belongs to a conjugate of
the level, then its characteristic polynomial has integral
coefficients (conjugate it into the integral adelic matrices).
For rank $h$, the algebraic integer $h-\operatorname{tr}u$ is
zero only if $u=1$: at a complex place the definite unitary matrix
is diagonalizable with eigenvalues on the unit circle, whose sum
can be $h$ only if they are all one. At all embeddings its modulus
is at most $2h$. The algebraic integer is divisible by every chosen place
above $p_0$, because the conjugated matrix reduces to the identity
there. The admissible congruence $p_0\equiv1$ modulo the cyclotomic conductor
makes $p_0$ completely split in $E$: its Frobenius fixes the roots
of unity generating the ambient cyclotomic field. The $2f$ primes
come in conjugate pairs. Since
$\pi_M\overline{\pi_M}=p_0p_1$ with $p_1\ne p_0$, precisely
one prime of each pair divides $\pi_M$, with multiplicity one.
Thus there are $f$ chosen primes of norm $p_0$. If the trace
difference were nonzero,
its absolute norm would be at least $p_0^f$, but at most
$(2h)^{2f}$. The rank bound $h\le2n_m$, with $n_m=k+m(5r+1)$, and
the level bound $p_0>(8n_m)^2$ make this impossible. The lifted inner product is the finite class sum of Weyl pairings.
Here is the local calculation for one term of this sum.
Let $F$ be the integral embedding of the Weyl module into its
column tensors with orthonormal filling basis. A row of $F$ is the
expansion of the product of minors for that filling in the standard
tableau basis. Standard fillings give linearly independent columns.
Cauchy--Binet gives $F S(a)=M(a)F$, where $S$ is the Schur action
and $M$ the tensor of the exterior-power matrices. Thus the positive
integral Gram matrix $G_W=F^tF$ satisfies
$G_WS(a)=S(a^t)^tG_W$, since $M(a^t)=M(a)^t$. Each determinant
character in the actual weight has the same value on $a^t$ and $a$,
so the relation holds for the full weight action $\pi$.

At a representative $g$, let $a$ be the tuple of its matrices for
$\iota$. The component of $g^{-1}$ for $\bar\iota$ is $a^t$, by
unitarity. For vectors integral at the opposed identifications,
$x=\pi(a^t)\bar\iota(u(g))$ and
$y=\pi(a^{-1})\iota(v(g))$ have integral coordinates by the gauge
definition. Their pairing $x^tG_W y$ equals
$\bar\iota(u(g))^tG_W\iota(v(g))$ by the transpose relation.
It is integral. The class weights are one, so summation gives
\begin{equation}\label{c3:local-pair}
 u\in\mathcal L_{\eta,\bar\iota},\ v\in\mathcal L_{\eta,\iota}
 \quad\Longrightarrow\quad
 |\langle u,v\rangle|_\iota\le1.
\end{equation}
The same class sum at the trivial twist has metric
$G_{\rm cl}=I_{\rm classes}\otimes G_W$, giving the squared conjugate norm
in \eqref{sz:covol}. We use the single matrix from \eqref{sz:index}
at each identification. For injectivity of $\Theta_\eta$ one applies
the sharp-datum injectivity assertion of input K to the elementary
arrows in the lower triangle. The positive level index and Weyl
dimension in the algebraic-to-Petersson norm conversion imply that
a zero algebraic lift has zero weighted chain. Its weights are
nonzero finite group orders, so reversing the injective arrows gives
zero in $W$.

\subsection{A signed height calculus}\label{c3:height}

For a family $u=(u_\iota)$ of complex scalars, write
$\operatorname{ht}(u)\le C$ if there are a multiplier $q\ne0$,
a real $\alpha$ and numbers $h_\ell\ge0$ such that
\begin{equation}\label{c3:heightdef}
 |\sigma q|\le e^\alpha,\qquad |qu_\iota|_\iota\le
 \ell^{h_\ell},\qquad
 \alpha+\sum_{\ell\in F}h_\ell\log\ell\le C.
\end{equation}
All automorphisms, identifications and finite prime sets are included.
This is a predicate, not a conventional nonnegative height function.
In particular $\alpha$ and $C$ may be negative. For products multiply
the multipliers and add both kinds of exponents. This proves the
additivity of the upper bounds, including powers with nonnegative
integer exponent. Multiplication by a family of powers
$N_{\eta_\iota}^{s_\iota}$ with rational exponents changes no finite
bound when the twist is allowed: the positive real root is an
$\iota$-adic unit because some integer power is a power of the
prime-to-$\ell$ integer $N_{\eta_\iota}$.

If $\operatorname{ht}(1)\le C$ then $C\ge0$. Indeed the multiplier
$q$ is algebraic: if it were transcendental, a field isomorphism
$\mathbb C\simeq\overline{\mathbb Q}_2$ could send it to a
transcendental element with arbitrarily large $2$-adic norm, contrary
to the uniform bound. To see the extension assertion, both fields
are algebraically closed of cardinality continuum and transcendence
degree continuum over $\mathbb Q$; extend the isomorphism of rational
function fields using transcendence bases and algebraic closures.
The set of elements algebraic over $\mathbb Q$ is countable, whereas
any nonempty $2$-adic ball contains uncountably many elements (already
in a translate of a ball in $\mathbb Q_2$). Thus the prescribed
transcendental elements exist. Set $K=\mathbb Q(q)$, $d_K=[K:\mathbb
Q]$, and $R=N_{K/\mathbb Q}(q)\ne0$. All embeddings of $K$ at a
finite or infinite place extend to identifications or automorphisms.
Consequently
\[
 \log|R|\le d_K\alpha,\qquad
 \log|R|_\ell\le d_K h_\ell\log\ell.
\]
The ordinary rational product formula, summed over the finite prime
set supporting the numerator and denominator of $R$, yields
$0\le d_K C$. In particular if a nonzero constant $z$ satisfies
$\operatorname{ht}(z)\le C$ and $|\sigma z|\le e^{C'}$
for every $\sigma$, then use multiplier $z$ to bound the height of
$z^{-1}$ by $C'$ and multiply to obtain
\begin{equation}\label{c3:product}
 0\le C+C'.
\end{equation}

\subsection{A regular family of finite-place twists}\label{c3:twists}

We need a single $m$-fold twist $\eta_\iota$ for each local
identification, simultaneously good for all $m$ single lifts, not
one twist for each member of $y$ or one for each $j$. No relation
between $\eta_\iota$ and $\eta_{\bar\iota}$ is required.
A critical pair $p=(\nu,z)$ is a unitary character with its negative
archimedean weight and tame part, and an integer $z$ in the critical
range specified by $D$; write
\[
 R_p(\vartheta,s)=
 \frac{|L^{S}(z,(\nu\vartheta^s)^{-1})|}
 {|L^{S}(z,\nu^{-1})|},\qquad s=\pm1,
\]
where in the denominator $S$ comprises the chosen places and their
conjugates, and in the numerator it additionally comprises $Q(\vartheta)$
and its conjugates. The change of omitted factors is included in the
explicit modulus powers below. For $s=1$ omit $s$.
The positive ratios are viewed as complex scalars, so their
conjugates need not have the same size. A ramified pair has tame part
nontrivial exactly at the chosen places. Let $J_j(\eta)=\eta^j$,
and use the superscript $j$ only for the power of the twist.

The required regularity is as follows:
\begin{enumerate}
\item $(N_{\eta_\iota},\ell)=1$;
\item for each ramified $p$, each $j$ and sign $s$,
$\operatorname{ht}(R_p(\eta_\iota^j,s)^{-1})\le x_1$;
\item for each $j$ there is a nonnegative logarithmic local budget
$\beta_{j,\ell}$ with $\sum_{\ell\in F}\beta_{j,\ell}\le B_0$.
Write $E_{j,\vartheta}$ for the image of the single lift on good
forms on $W_j$ and $P_{j,\vartheta}$ for its orthogonal projection.
There is an operator $\Lambda_\iota$ on $E_{j,\eta_\iota^j}$
with $|\det\Lambda_\iota|_\iota\ge
\exp(-d_T r\beta_{j,\ell})$ such that
$\Lambda_\iota P_{j,\eta_\iota^j}$ takes integral forms
orthogonal to false lifts to integral forms there. Also, for each
main critical pair $p$ of lift $j$, if $\Xi,\Xi'$ are test forms
for its character lift $\phi_\iota$ (integral for every embedding into $\mathbb C_\ell$ and
orthogonal to false lifts in the kernel of the adjoint), then
\[
 |\langle P_{\phi_\iota}\Xi,P_{\phi_\iota}\Xi'\rangle|_\iota
       \le e^{\beta_{j,\ell}}.
\]
Here $P_\phi$ is projection onto $\mathbb C\phi$.
\end{enumerate}

Here is the reduction of regularity to the local arithmetic input.
At fixed $\ell,\iota$, the good Hecke eigensystems on $W_m$ form a
finite set: their commuting normal Hecke operators simultaneously
diagonalize a finite dimensional space. For every $j$ and every such
system one takes the rank-$r$ main summand of its Galois representation;
for each main critical pair take the rank-one summand for the trivial
system. Join both finite collections. The simultaneous twisting
theorem for these collections and the ramified pairs produces $\eta$
whose first $m$ powers have auxiliary modulus away from $\ell$,
short divisible Selmer groups and inverse-ratio norms at most
$\ell^{u_\ell}$, where $G_0$ is the product of the orders of $\mathrm{GL}_{rt}(k_w)$ over the chosen places above $p_0,p_1$, with $t=5r$, and
\[
 u_\ell=(4f+3)r v_\ell(G_0)+2^{-\ell}.
\]
In the source the construction starts at a chosen place away from $\ell$
using its tame inertia generator, constructing lattices in the main
summands and Stickelberger systems for the ramified ratios, and
applying the simultaneous anticyclotomic twisting theorem to all
powers with error $2^{-\ell}$. Identification of the resulting avatar
with an automorphic tame character gives the same modulus for every
power: the avatar is nontrivial at each auxiliary place for each
power. Appendix~\ref{sec:finite} passes from the simultaneous polynomial
datum to a single twist using one root-of-unity tuple for all powers.
The Selmer and Stickelberger construction of the datum remains in R.

The first-order congruence denominator argument on the span of
lift generators applies these Selmer lengths to both the rank-$r$
space and the rank-one line. In more detail, choose a basis $v_i$ from the lift generators.
The first-order congruence calculation gives one operator $G_i$ for
each generator, diagonal on this basis with its $i$-th diagonal
entry of valuation at most the lift-valuation bound. Each $G_i$
after projection preserves the integral module on forms orthogonal
to false lifts. As the residue field of $\overline{\mathbb Q}_\ell$
is infinite, a combination with integral scalar coefficients
can avoid cancellation in every diagonal position: normalize the
entries in each row by one of largest norm, and avoid the finitely
many proper linear hyperplanes defined by their reductions. The
resulting diagonal operator has every entry of valuation at most
the stated bound, and retains integrality by linearity. Multiplying
by $G_0$ clears all stabilizer denominators for pairings. Taking
the product of the eigenvalues gives determinant bound
$\ell^{-\dim(E)r\,\mathrm{clearingBound}(D,\ell)}$ (rank one for
the line). In the rank-one case the denominator bound and the
pairing bound therefore control the line projection itself, with
the class-number factor in the local budget accounting for the
normalization of test forms. Appendix~\ref{c4:central-average} shows that the rank-one
central average changes the pairing by exactly the class factor
$N_0$. Existence of the generator-wise operators remains in R.
Thus one may take
\[
 \beta_{j,\ell}=\bigl(\mathrm{clearingBound}(D,\ell)
                     +v_\ell(N_0)\bigr)\log\ell ,
\]
where $N_0$ is the number of $U(1)$ classes at chosen level. This
choice is actually independent of $j$.

Since $\sum_{\ell\in F}v_\ell(G_0)\log\ell
\le\log G_0$ and
$\sum_{n\ge2}2^{-n}\log n\le\sum_{n\ge2}n2^{-n}<2$,
\[
 \sum_{\ell\in F}u_\ell\log\ell
 \le (4f+3)r\log G_0+2\le x^{999}\le x_1 .
\]
The clearing bound is the sum of the local defect, $(1+2\delta)u_\ell$,
the denominator error, and $v_\ell(G_0)$. Applying the defect envelope and the numerical estimates
in Appendix~\ref{sec:finite} gives
\begin{align*}
 \sum_{\ell\in F}\beta_{j,\ell}
 &\le 4\cdot10^6e^{180} f(m/k)bL
 +(1+2\delta)x^{999}+x^{4000}+\log G_0+x^{30}\\
 &\le B_0.
\end{align*}
The denominator error is a linear combination of prime valuations of
the fixed odd class bound and $G_0$; the source bounds its sum by
$x^{4000}$ and the $N_0$ term by $x^{30}$. Select one such twist for
each identification. With multiplier $1$ the local inverse-ratio
norms imply item 2 by (\ref{c3:heightdef}). The remaining local arithmetic hypotheses in this reduction are
specified under R in \S\ref{c3:frontier}.

\subsection{The height gain in a single lift}\label{c3:single}

Fix $j\in[1,m]$. Let $T_\vartheta$ denote the stage-$j$ theta lift with twist $\vartheta$ restricted to the good forms of the rank-$r$ hermitian space $W_j$ of signature $(k+j-1,m-j)$. All these
sources have dimension $d_T$ by the identity-operator case of the
trace comparison in Appendix~\ref{sec:trace}. Define $Z_j(\eta)=\det(T_{\eta^j}^*T_{\eta^j})$.
The comparison operator is
\[
 U_\vartheta=G_0 h_j^2 g_\vartheta g_\vartheta^* ,
\]
where $g_\vartheta$ is the gluing map from tuples of forms on the
smaller definite space $V^\circ$ and $h_j$ is the hook number of its
negative Weyl shape. Set $A_{\rm gl}=fmbL^2$. Input G gives the seesaw pairing with a tuple $\zeta$ as the sum
over kernel choices $i$ of fixed nonzero $c_i$ times the integral
of $y(g)\,F_i(g)\,\theta^\circ(g,\zeta_i)$, where $F_i$ is the
carrier factor and $\theta^\circ$ the complementary block kernel.
This formula is independent of the twist after multiplication
by $c'_\vartheta$. Testing against tuples supported in one block
shows that a vector killed by the resulting map has zero pairing
with each of these products. The products-spanning input says
precisely that such a good vector is zero. In particular the map
$c_\vartheta' g_\vartheta^*T_\vartheta=\Pi_j$
is independent of $\vartheta$, with $\Pi_j$ injective, and
the quotient $\rho_\vartheta$ of the twisted to untwisted product of
Hecke values satisfies
$\rho_\vartheta=|c_1'|^2/|c_\vartheta'|^2\ne0$.
Hence, for the source inner product,
\[
 T_\vartheta^*U_\vartheta T_\vartheta
 =\rho_\vartheta U^\circ,
 \qquad U^\circ=G_0h_j^2|c_1'|^{-2}\Pi_j^*\Pi_j>0.
\]
Writing $C_\vartheta$ for the compression
$P_{j,\vartheta}U_\vartheta|_{E_{j,\vartheta}}$ and
$F_\vartheta=\det C_\vartheta$, we have
\[
 Z_j(\eta) F_{\eta^j}=\det(U^\circ)\rho_{\eta^j}^{d_T}.
\]
Indeed relative to an orthonormal basis of the image,
$T^*UT=T^*C T$, so determinants multiply, and $\det T^*T=|\det T|^2$.
All factors are nonzero by positivity of $U^\circ$ and injectivity.
Division by the untwisted identity gives
\begin{equation}\label{c3:compression}
 \frac{Z_j(1)}{Z_j(\eta)}
 =\frac{F_{\eta^j}}{F_1}\rho_{\eta^j}^{-d_T}.
\end{equation}

\paragraph{Integral comparison and conjugate norms.}
In the tuple space take the integral tight frame $(f_a)$, at rational
translates of the class representatives integral above $\ell$.
Its frame operator is $h_j^2 I$; conjugated coordinates in the
translated Weyl lattice are integers. The gluing integrality hypothesis takes this
frame to forms integral at both $\iota$ and $\bar\iota$ away from
the modulus. The frame identity gives
\[
 U_\vartheta v=\sum_a G_0
       \langle g_\vartheta f_a,v\rangle\,g_\vartheta f_a .
\]
The factors $G_0$ clear the orders of the class stabilizers, so the
coefficients are integral for integral $v$ by the same local Weyl
pairing as (\ref{c3:local-pair}). The output is integral.
The gluing-period identity also implies that $U_\vartheta$ takes
lifts to forms orthogonal to false lifts.

To construct the frame, let $h=h_j$ and use the filling matrix $F$
just described for the negative shape. Write $G_W=F^tF$.
Let $B_{ij}=b_i(\partial)b_j(0)$ be the integral Fock pairing of
the standard minor polynomials. Its normalization gives
$G_WB=h I$. Indeed substitution by $a$ and coefficient differentiation
are adjoint under transposition of $a$, so $B^tG_W$ commutes with
the Schur representation at every place. The tensor of its
irreducible highest weight modules is irreducible for the product
of the general linear groups. Schur's lemma therefore makes the
operator scalar. Evaluation at the highest polynomial, whose
differentiation pairing with itself is $h$ and whose kernel value
on the standard flag is $1$, fixes that scalar. The last
normalization is the definition of the hook number in the Fock
model. For a tensor basis vector $e_T$ take $a_T=BF^te_T$.
Then $a_T$ is integral and $G_Wa_T=hF^te_T$. Hence
$Fa_T=h\,\operatorname{pr}_{\operatorname{im}F}e_T$.
The projections of an orthonormal basis are a Parseval frame for
the image. Multiplication by $h$ gives frame constant $h^2$ and
$\|a_T\|^2\le h^2$.

If $a_T$ is a frame vector
and $u$ is a rational unitary translate with representation
$\pi(u)$, replace $a_T$ by
$\overline{\pi(u)^{-1}\overline{a_T}}$.
This is a unitary change of frame: the frame operator remains
$h_j^2 I$. Moving the conjugate coordinates back gives the
original integer vector. The Gram identity for the unitary
representation is rational in the coordinates and their
conjugates over the CM field, and remains valid at each
archimedean embedding. Since the original vector is integral,
the moved conjugate norms retain the bound $h_j^2$.
Independently place these vectors in each block of the finite
tuple space, with zero other blocks; summing the frame identity
blockwise gives the tuple tight frame.

At the trivial twist use the tight frame at the unshifted
representatives. Its conjugate squared tuple norms are at most
$h_j^2$. Let $B$ be the set of tuple-coordinate blocks, $I$ the
frame indices; the gluing and counting estimates are
\[
 \|g_1 v\|_\sigma^2\le e^{2A_{\rm gl}}\|v\|_{\sigma,\mathrm{tuple}}^2,\qquad
 |B|(|I|G_0 h_j^2)^2\le e^{2A_{\rm gl}}.
\]
For the cardinality estimate put $N=|B|$ and let $C_f$ be the
number of column fillings for the negative shape. The frame has
$|I|=NC_f$. The class and hook estimates of input G give
$(NG_0)^2\le e^{A_{\rm gl}}$, $C_f,h_j\le e^{x^{999}}$. Since
$N\le e^{A_{\rm gl}/2}$ and $6x^{999}\le A_{\rm gl}/2$,
\[
 N(NC_fG_0h_j^2)^2\le
    \exp(A_{\rm gl}/2+A_{\rm gl}+6x^{999})\le e^{2A_{\rm gl}}.
\]
For $S=e^{2A_{\rm gl}}\|v\|_\sigma^2$, the conjugate coefficient of the
$a$-th frame member is at most $G_0h_j\sqrt S$ by Cauchy--Schwarz.
(The general automorphism changes conjugation to its conjugated
automorphism; the frame bound holds for both.) By triangle inequality
each block of the resulting tuple has norm at most
$|I|G_0 h_j^2\sqrt S$. Summing the squares over $B$ and gluing once
more bounds $\|U_1 v\|_\sigma^2$ by $e^{6A_{\rm gl}}\|v\|_\sigma^2$.
The projection on lifts, applied to a vector orthogonal to the false
lifts, is contractive for every conjugate norm; this uses the
algebraic characterization by the central idempotent and the
Hecke polynomial kernel, not continuity of $\sigma$.
Consequently, for an eigenvector $v\ne0$ of $C_1$ with eigenvalue
$\lambda$, $|\sigma\lambda|^2\|v\|_\sigma^2\le
e^{6A_{\rm gl}}\|v\|_\sigma^2$. The determinant is the product of roots
of the characteristic polynomial with multiplicities, yielding
\[
 |\sigma F_1|\le e^{3d_T A_{\rm gl}}.
\]

At a local identification use $\vartheta=\eta_\iota^j$ and
the clearing operator from regularity. Then $\Lambda_\iota
C_\vartheta$ preserves the integral module of its image. An
eigenvector can be scaled into that module with at least one
transported coordinate of norm one: take the finitely many class
values at representatives integral above $\ell$, take the maximum
of the coordinate norms in the local Weyl model and scale by the
inverse of a coordinate attaining it. Every eigenvalue is therefore
integral. By the product of eigenvalues and the bound on
$\det\Lambda_\iota$,
\[
 |F_{\eta_\iota^j}|_\iota\le e^{d_T r\beta_{j,\ell}}.
\]
In (\ref{c3:heightdef}) take multiplier $F_1$, archimedean budget
$3d_T A_{\rm gl}$, and finite budgets $d_T r\beta_{j,\ell}/\log\ell$.
This proves
\begin{equation}\label{c3:compressionheight}
 \operatorname{ht}(F_{\eta_\iota^j}/F_1)
 \le d_T(3A_{\rm gl}+rB_0).
\end{equation}

\paragraph{The main rank-one ratios.}
The $-A_0$ gain comes from the following two different sides
of a rank-one computation. Put
$J_0=\sum_\tau j(\kappa_\tau)$,
$S_{\Sigma}=\sum_\tau S(\kappa_\tau,\tau)$ so that
$A_0=J_0-S_{\Sigma}$. For a main pair $p$ of lift $j$, let
$\phi_\vartheta$ be its normalized rank-one lift to $U(2z)$
and $\phi_1$ the untwisted lift.

To bound the constant $\|\phi_1\|^2$, compare it to the lift of a
flat pair $q$ of the same rank and empty row. The integral algebraic
part $I_q$ of the flat value at trivial twist is integral everywhere
and equals $G_0\mathrm{denomBound}(q)^2\|\phi_{q,1}\|^2$.
For automorphisms fixing the base field, the single-lift norm formula
gives
\[
 \log|\sigma\|\phi_{p,1}\|^2|
 =C_\sigma+\epsilon_{p,\sigma}
    +\sum_\tau j(l_\tau),\qquad
 l_\tau=p.\mathrm{negWeight}(\tau)-z\ge\kappa_\tau .
\]
Here the constant is common to the two pairs,
$\epsilon_{p,\sigma}\ge-4f$ and
$\epsilon_{q,\sigma}\le4f$ with empty row $l_\tau=0$.
The same moduli of conjugates are realized by automorphisms fixing
the base field. Thus for the multiplier $I_q/\|\phi_1\|^2$
the common constant cancels, the finite budgets are zero and the
archimedean budget is at most
$x^{33}+8f-J_0\le B_0-J_0$. Therefore
\begin{equation}\label{c3:main-arch}
 \operatorname{ht}(\|\phi_1\|^2)\le B_0-J_0 .
\end{equation}
Here one uses monotonicity of $j$ on the
range in question, and bounds the fixed class and denominator
factors by $e^{x^{33}}$. These and the common-constant norm formula
are the numerical hypotheses in input F.

For the finite side let $s_\ell$ be the sum of the local gains at $\ell$ over $\tau\in\Phi$, computed below.
The crucial sharpness estimate is existence of a common nonzero scalar
$c_0$ with $|\sigma c_0|\le e^{B_0}$ and, for every allowed twist and
identification, a test form $\Xi$ satisfying
\begin{equation}\label{c3:test}
 |\langle\Xi,\phi_\vartheta\rangle|_\iota
 \ge |c_0 R_{q_2}(\vartheta)|_\iota \ell^{-s_\ell},
 \qquad 2\sum_{\ell\in F}s_\ell\log\ell\le S_{\Sigma} .
\end{equation}
The prime-sum inequality is just regrouping by prime powers.
In fact put
\[
 g_\tau(v)=(1-\chi_K(v))
       (\kappa_\tau/v-\kappa_{-v\tau})_+ .
\]
The notation $-v\tau$ uses the place of the class of $-v$ if
prime to the cyclotomic conductor and the identity place otherwise.
The weight is at least $a/x_1$, so $g_\tau(v)$ vanishes for
$v\ge x_1$; all terms are nonnegative. To get the local formula, compare the coefficient of each positive
exponent $e$ with the defining sum for the local gain.
At a split prime $1-\chi_K(\ell)^e=0$. At an inert prime this
coefficient is $2$ for odd $e$ and zero for even $e$, and the
local uniformizer weight is $1$. At the ramified prime the
coefficient is $1$ for every $e$, and the local uniformizer
weight is $1/2$. Thus
$2s_\ell=\sum_\tau\sum_{e\ge1}g_\tau(\ell^e)$.
Since $\Lambda(\ell^e)=\log\ell$ and the von Mangoldt
function vanishes at non-prime-powers, summing this equality over
the primes gives exactly $S_{\Sigma}$, and restriction to a finite set
only decreases it. Thus the prime-sum regrouping consumes no error.

Here $q_2$ is an empty-row ramified pair of rank four with the fixed
chosen-level tame character. To specify the origin of the sharper
Taylor input:
the rank $2z$ indices split into blocks of sizes $12,4,4,2z-20$.
The first block gives a Taylor coefficient of weight
$l_\tau\in[\kappa_\tau,\kappa_\tau+x^{42}]$ at a finite indicator
datum; the next carries $q_2$, and the last two carry empty rows
with tame characters whose product is the character selected by
class averaging. The sharp Taylor coefficient is scaled by
$T=M^{12f}c_{\rm vac}\prod_\tau
\mathrm{lossProduct}(l_\tau)^2$ with
$|\sigma T|\le e^{fmb/k}$. Its local valuation upper bound before
scaling is $s_\ell+2\sum_\tau v_\ell(\mathrm{lossProduct}(l_\tau))
-\sum_\tau v_\ell(l_\tau!)$ relative to the integral vacuum.
Multiplication by the factorials and division by $T$ removes both
loss terms, leaving $s_\ell$. The scaled coefficient is $C s(1)$
for a one-variable kernel $s$ on $U(1)$.

Appendix~\ref{jets:sharp} separates the valuation lower bounds
for all coefficients from the upper valuation bound supplied by
one coefficient. The latter direction is needed here.

Class averaging can be made explicit. The rank-one definite
class space is a finite abelian quotient $C$ of the group of finite adelic
points modulo rational points and the chosen level. There
are no stabilizers. Translate the chosen representatives by the
inverse of the representative of the identity class, so that
this representative becomes $1$. The map taking the finite component of a norm-one adele to its
rank-one matrix sends rational points into the rational subgroup,
and integral ideles reducing to $1$ at the chosen places into the
level. Pulling a character of $C$ back along this map therefore
gives a weight-zero tame character with possible conductor only at
the chosen places. At a finite rank-one group element its value is
the original character, since this map composed with the rank-one
adelic determinant is the identity. For a character $\chi$ form the
normalized mean $A_\chi=|C|^{-1}\sum_{g\in C}
\chi(g)s(g_{\rm rep})$. Character orthogonality gives
$\sum_\chi A_\chi=s(1)$; the ultrametric inequality therefore
gives a character with $|A_\chi|_\iota\ge|s(1)|_\iota$.
Call this mean $A_s$.
For the other blocks let $W_\psi$ be the product of the two flat
integral algebraic parts with characters multiplying to the selected
character. The geometric mean $I$ of all such products can be
taken independent of the selected character because each
factorization runs over the entire finite character group. By
multiplicativity some factorization has $|W_\psi|_\iota\ge|I|_\iota$.
All conjugates of $I$ have norm at most $e^{2x^{34}}$ by the
corresponding flat bounds. The kernel pairing identity is
\[
  T I_{q_2} W_\psi(C A_s)
    = N_0^2\langle\Xi,\phi_\vartheta\rangle/R_{q_2}(\vartheta).
\]
Taking norms, using $|N_0|_\iota\le1$, proves~(\ref{c3:test})
with $c_0=T I_{q_2} I$; its archimedean bound is
$e^{fmb/k+3x^{34}}\le e^{B_0}$.
The Taylor sharpness bound, existence of the block test form with
the displayed identity, and the flat integral algebraic estimates
are the substantive inputs of this argument (input F below).

The precise external result here is the
\srcline{New/OneLift/HeckeValues/TwistedValueHeights}{700}{sharp test-form theorem}.
For fixed main pair and rank-four companion, it chooses one nonzero \(c_0\)
before the prime, the identification and the allowed twist. In its proof
the degrees are prescribed by the negative weights and the fixed shift.
The \srcline{New/Derivatives/Sharpness/BoundAttained}{381}{attainment theorem}
supplies the coefficient data; the geometric mean \(I\) used for \(c_0\)
is also fixed before local specialization. The indicator, averaging
character and factorization yielding \(W_\psi\) may depend on the
specialization. The
\srcline{New/OneLift/HeckeValues/TwistedValueHeights}{571}{integral kernel-pairing theorem}
realizes these choices on the required test form. This order of choice
permits the same \(c_0\) and twist at two identifications.

Now take $\vartheta=\eta_\iota^j$, $\phi=\phi_\vartheta$,
$\rho=R_{q_2}(\vartheta)$ and choose another test form
$\Xi'$ at $\bar\iota$ for the \emph{same} $\vartheta$.
Since $\rho$ is real, multiplying the two inequalities gives
\[
 |\langle\Xi,\phi\rangle\langle\phi,\Xi'\rangle|_\iota
 \ge |c_0\bar c_0\rho^2|_\iota\ell^{-2s_\ell}.
\]
Denote the nonzero pairing product by $I_\iota$. Direct projection
onto the line gives
$\langle P_\phi\Xi,P_\phi\Xi'\rangle=I_\iota/\|\phi\|^2$.
By the rank-one norm formula for the Hecke ratio,
\begin{align}
 R_p(\vartheta)^{-1}
 &=N_\eta^{2z-1}\frac{\|\phi_1\|^2}{\|\phi\|^2}\notag\\
 &=N_\eta^{2z-1}\|\phi_1\|^2
 \langle P_\phi\Xi,P_\phi\Xi'\rangle\rho^{-2}
 \frac{\rho^2}{I_\iota}.\label{c3:four-factors}
\end{align}
The first factor after the modulus power costs $B_0-J_0$ by
(\ref{c3:main-arch}); the projection costs $B_0$ by regularity;
the ramified ratio costs $2x_1$. For the last factor use multiplier
$c_0\bar c_0$, archimedean budget $2B_0$ and finite budgets
$2s_\ell$; equivalently the local norm bound is
$\ell^{2s_\ell}$. It costs at most $2B_0+S_{\Sigma}$.
Since $x_1\le B_0$ and the modulus power costs nothing,
\begin{equation}\label{c3:main-height}
 \operatorname{ht}(R_p(\eta_\iota^j)^{-1})
 \le 6B_0-A_0 .
\end{equation}


\paragraph{Counting and combining the three exchanges.}
The three exchange products for the comparison operator are
products of squared critical values with modulus powers.
The unsigned exchange character contributes points $1\le z\le7(m/2)+k/2$,
the two signed flat characters points $1\le z\le m$ and
$1\le z\le m/2$. Reserve for the main-pair estimate the points
$7(m/2)<z\le7(m/2)+k/2$. Since $m$ and $k$ are even, there
are $k/2$ of these, and $7(m/2)+m+m/2=5m$ remaining
ramified points. The ratio of exchange products is therefore
$N_\eta^e(\prod_{a=1}^{k/2}R_{p_a}\prod_{b=1}^{5m}
R_{q_b})^2$, with the prescribed signs on the second product.
This is the partial $L$-value factorization of input G; taking
inverse heights, the squared product doubles the counts, giving
\[
 \operatorname{ht}(\rho_{\eta_\iota^j}^{-1})
 \le k(6B_0-A_0)+10mx_1.
\]
The scale inequality
$3A_{\rm gl}+rB_0+10mx_1+6kB_0\le kLB_0$
and (\ref{c3:compression}), (\ref{c3:compressionheight}) now prove
\begin{equation}\label{c3:single-height}
 \operatorname{ht}\left(\frac{Z_j(1)}{Z_j(\eta_\iota)}\right)
 \le d_T\,k(LB_0-A_0)=d_T\left(\frac{V}{m}-kA_0\right).
\end{equation}


\subsection{From single lifts to the chain}\label{c3:chain}

Put $Z(\eta)=\det(\Theta_\eta^*\Theta_\eta)>0$.
One cannot simply multiply the norms of the single lifts, since
they have different sources; the triangle-of-lifts norm identity
and the equality of Hecke multiplicities are needed. In the source
the $j$-th single-lift Gram operator is a constant $a_j\ne0$
times the common local twist factor $c_Q$ times the Rallis operator
$\mathcal R_j(\eta)$ on its good space. For the unnormalized coset sum write $R$ for the operator with
matrix coefficients $\sum_\gamma K_\eta(\gamma)
\langle x,\gamma y\rangle$; $K_\eta$ is the unramified vacuum
coefficient. For the sharp finite block datum of mass $M_{\rm blk}$
the asserted Petersson Rallis formula of input K (cf.~\cite{gqt} for the general inner product formula) gives the
same sum times $C_jM_{\rm blk}$ with $C_j$ depending only on
the archimedean datum. The latter is fixed under tame weight-zero
twists. Converting to algebraic forms multiplies by the index
of the level and the Weyl dimension. The finite mass relation is
$[K_0:J]M_{\rm blk}=\operatorname{gramFactor}/M_Q$;
the quotient is the chosen-place factor times $c_Q$.
Thus the scalar multiplying $R$ is $a_j c_Q$, where $a_j$
is nonzero and independent of the twist. The paired-series
argument of Appendix~\ref{sec:trace} realizes the two $R$'s
on $W_j$ and $Y_0$ by one Hecke operator and identifies its
two determinants.
For specificity, write $w_n$ for the order of the general linear
group in the $n$-th step of the lower triangle and $M_Q$ for the
number of $Y_1\times\mathrm{RowIndex}$ matrices over the residue
ring of the chosen and twist modulus. The first-arrow Gram formula
in row $i$ is $w_0^{-1}C_i M_Q^{-1}\mathcal R_{i+1}$, transported
to the common starting space, with $C_i$ independent of twist.
Here the Gram pairings are sesquilinear.
The weighted triangular computation works by induction on the row
length $n=m-i$. In the zero row the weighted chain is the identity.
Moving the adjoint of the first step past the $n-1$ subsequent
steps multiplies its scalar by the product of the commutation
constants in those squares. Their moduli, up to a fixed norm
power, are $Q_{i+1,j}$, $i+1<j\le m$; the adjoint law contributes
only scalar factors of modulus one in this sense. Replacing the
paired first-step ascent and descent by its first-arrow Gram
formula inserts $C_i M_Q^{-1}\mathcal R_{i+1}$.
The weights telescope from $w_0w_1^{n-1}/w_{n-1}$
to $w_0w_1^n/w_n$. More explicitly the weighted row transfer
gives the factor $w_0 w_1 w_{n-1}/w_n$ and the first-arrow
pairing contributes $1/w_0$.
Thus at row zero the weighted chain $u_\eta$ has Gram pairing
\[
 \langle u_\eta v,u_\eta v'\rangle
 =\kappa_\eta\frac{w_0w_1^m}{w_m M_Q^m}
    \langle v,\prod_{j=1}^m\mathcal R_j(\eta)v'\rangle,
\]
where the modulus of $\kappa_\eta$ contains exactly one factor
$Q_{ij}$ for each $1\le i<j\le m$. The induction uses the square
commutation and adjoint identities; their arithmetic derivation
is part of input K, not an identity of arbitrary theta lifts.

Conversion from the weighted Petersson model to the final
algebraic model multiplies this Gram form by the Weyl dimension
$h_Y$ and the index of the final level. The precise index
computation is
\[
 \mathrm{ind}(J_\eta)=C_1 N_\eta^u
     \mathrm{gramFactor}(\eta)^m\frac{w_m}{w_0w_1^m}.
\]
The integer $u$ and constant $C_1$ are independent of $\eta$:
at the chosen places all group orders are fixed, and at twist
places the block upper-triangular parabolic has order equal to the
product of the general linear orders of diagonal blocks times
$q_w^{\sum_{a<b}h_a h_b}$ for block sizes $h_a$.
This is obtained by independently choosing each upper off-diagonal
matrix block. Also
\[
 \frac{\mathrm{gramFactor}(\eta)}{M_Q}
 = \mathrm{chosenFactor}\,c_Q,\qquad
 c_Q=N_\eta^{r\delta}\prod_{w\in Q(\eta)}
                      \prod_{s=\delta+1}^{\delta+r}(1-q_w^{-s}).
\]
The product of $1-q_w^{-s}$ is the proportion of rank-$r$
matrices among $(r+\delta)\times r$ matrices: choose successive
independent columns in $q_w^{r+\delta}-q_w^i$ ways for $0\le i<r$.
The precise matching of the parabolic levels and of these factors
is input K. Multiplying the displayed formulas cancels every
$w_n$ and $M_Q^m$, leaving $h_Y C_1 N_\eta^u
\kappa_\eta\mathrm{chosenFactor}^m$ as the scalar multiplying
$c_Q^m$. The resulting operator identity is
\[
 \Theta_\eta^*\Theta_\eta
 =c_\eta c_Q^m\prod_{j=1}^m\mathcal R_j(\eta),\qquad
 |c_\eta|=C N_\eta^{e/2}\prod_{1\le i<j\le m}Q_{ij}(\eta),
\]
where $C>0,e\in\mathbb Z$ are independent of $\eta$. Here
$Q_{ij}$ is the product at $1\le s\le\delta$ of critical ratios
for the character $\nu_j/\nu_i$ at twist $\eta^{j-i}$,
where $\nu_j$ denotes the untwisted character at stage $j$.
This character is a power of the common ratio between consecutive
lift characters and is ramified at the chosen places.
The local triangle identities and the analytic Rallis formula
enter through K. The trace transfer uses the separate hypotheses
of Appendix~\ref{sec:trace}. Once they hold, take
determinants on $W$:
\[
 Z(\eta)=
 \left(c_\eta/\prod_j a_j\right)^{d_T}\prod_j Z_j(\eta).
\]
The common $c_Q^m$ and all Rallis determinants cancel. Since the
Gram determinants are positive, take absolute values, and divide by
the trivial-twist identity (all $Q_{ij}(1)=1$):
\begin{equation}\label{c3:chainfactor}
 \frac{Z(1)}{Z(\eta)}=
 \left(N_\eta^{e/2}\prod_{i<j}Q_{ij}(\eta)\right)^{-d_T}
 \prod_{j=1}^{m}\frac{Z_j(1)}{Z_j(\eta)} .
\end{equation}
Each $Q_{ij}$ contains $\delta$ ramified ratios, and $1\le j-i\le m$,
so their inverse heights cost at most $m^2\delta x_1\le b$ before
raising to the power $d_T$. Summing (\ref{c3:single-height}) over
the same $m$ twists $\eta_\iota^j$, using $mkA_0=H$, gives
\begin{equation}\label{c3:chainheight}
 \operatorname{ht}\left(\frac{Z(1)}{Z(\eta_\iota)}\right)
 \le d_T(b+V-H).
\end{equation}

\subsection{Short auxiliary completion}\label{c3:auxiliary}

We require an auxiliary spanning set in $W$, independent of $y$.
The construction produces a scalar $\gamma(\eta)\ne0$,
$\gamma(1)=1$, with
\begin{equation}\label{c3:gamma}
 |\gamma(\eta)|=N_\eta^{e'/2}
       \prod_{i,j=0}^{m-1}T_{ij}(\eta),
\end{equation}
where $e'\in\mathbb Z$ is fixed, and $T_{ij}(\eta)$ is the
product of the $\delta$ ratios for the character $\nu_{i+1}\nu_{j+1}$
at $\eta^{j+1}$. Its tame character is nontrivial at the chosen
places, so each is a ramified-ratio product. The powers now use
$j+1$, not $j-i$. The spanning vectors $z$ have
\begin{equation}\label{c3:auxbounds}
 \gamma(\eta)\Theta_\eta z\in\mathcal L_{\eta,\iota}
 \ ((N_\eta,\ell)=1),\qquad
 \|\Theta_1z\|_\sigma^2\le\exp(2x_1^2b^2).
\end{equation}

Here is the auxiliary construction reduced to its path-period
identities. For the final grid-level space apply the short-basis estimate
of \S\ref{A:source} with the conductor integer in the different.
The grid's principal level is contained in the chosen level and its
rank $h_0$ is at most $2n_m$, giving trivial stabilizers by the
argument above. The weight has at most $b r^4$ boxes. With
$D_0$ the conductor integer, the numerical specialization is
$6(h_0 f+b r^4)\log(6D_0(fh_0)^2)\le x^{99}b^2$.
Transport through the algebraic realization at the last grid vertex
gives spanning forms all of whose conjugate values have squared
Weyl norm at most $e^{x^{99}b^2}$. There is a fixed onto map $R$
from this space to principal forms on $Y_0$. Pairings at the end
of the chain against a dual Weyl vector $v$ have the expression
\[
  v^t M_{\rm dual}
    \frac{c_\eta}{c_1}\Theta_\eta R\varphi(g)
    =\beta\,\mathcal P_\eta(\varphi,g,v).
\]
Here $\Theta$ in this identity denotes the chain on all principal
forms, $\mathcal P_\eta$ the last-path period (with one factor
a Weyl-valued delta form), and $|c_\eta/c_1|$ is the right side
of (\ref{c3:gamma}). The path data arise by carrying the first
path through all $m^2$ turns; their scalar moduli multiply. The
fixed dual hook number $h>0$ controls recovery from pairings.
There is a single integer $N\ne0$, of size at most
$e^{x^{99}b^2}$, clearing the period denominators for all allowed
twists on integral inputs. The period bound at trivial twist is
\[
 |\sigma\mathcal P_1(\varphi,g,v)|
 \le e^{x^{200}b^2}\|\sigma v\|_{\rm Weyl}
                  \sqrt{B}
\]
when every conjugate value of $\varphi$ has squared norm at most
$B$. These are the path-period and dual-model inputs Aux in
\S\ref{c3:frontier}. Rescale all inputs by a fixed $s\ne0$:
if $\beta\ne0$, take $s=hN/\beta$, else $s=1$.
The rescaled pairings are $hN\mathcal P_\eta$ or zero. The normalized dual model in input Aux has the following two properties. The integral dual-pairing matrix $M_{\rm dual}$ admits an integer
matrix $C$ with $CM_{\rm dual}=hI$ in the normalized dual model.
Testing the pairing against the $i$th row of $C$ gives
$h\zeta_i\in h\mathcal O_\iota$ for the lifted value $\zeta$.
Since $h\ne0$, every coordinate is integral, even at primes
dividing $h$. At infinity the dual model represents
$v\mapsto v^tM_{\rm dual}\sigma\zeta$ by a vector $w$ with
$\|w\|=\|\sigma\zeta\|$. Evaluate the bound at
$v=\sigma^{-1}w$ before conjugation. The resulting inequality is
$\|w\|^2\le hN e^{x^{200}b^2}\sqrt B\,\|w\|$.
The hook estimate $hN e^{x^{200}b^2}\le e^{x^{201}b^2}$ gives
squared Weyl norm at most $e^{2x^{201}b^2}B$. With
$B=e^{x^{99}b^2}$ this is at most $e^{x^{202}b^2}$.
The images $sR\varphi$ span all principal forms since $R$ is onto.
Apply the good projection multiplied by its fixed integral scalar
at the identity translate; it preserves integrality on the lifts,
is surjective onto $W$ and increases the logarithm of the conjugate
squared norm to at most $x^{99}+x^{202}b^2\le2x_1^2b^2$.
This yields (\ref{c3:auxbounds}) with $\gamma=c_\eta/c_1$.

By (\ref{c3:gamma}) and regularity we have, for the same
$\eta_\iota$ as in (\ref{c3:chainheight}),
\begin{equation}\label{c3:gammaheight}
 \operatorname{ht}(|\gamma(\eta_\iota)|^{-2})
 \le 2m^2\delta x_1\le b .
\end{equation}
The nonzero modulus power has zero finite cost; the coefficient 2
here comes from squaring the inverse absolute value, independently
of the index correction.

\subsection{The determinant comparison}\label{c3:compare}

\noindent\textbf{Theorem C (conditional on the geometric and arithmetic inputs
in \S\ref{c3:frontier}).}
Let $y_1,\ldots,y_n$ be independent in $W$,
$n\ge(1-4/b^2)d_T$, with $\operatorname{IndexLe}(y,\Delta)$
and $\operatorname{CovolumeLe}(y,c)$. Then
\begin{equation}\label{c3:theorem}
 d_T(H-3V)\le c+2\Delta.
\end{equation}

\begin{proof}
The case $d_T=0$ follows from $n=0$, the empty Gram determinant
($c\ge0$) and the empty prime-set index bound ($\Delta\ge0$).
Suppose $d_T>0$. Extend the given ordered family to a basis
$w=(y_1,\ldots,y_n,z_1,\ldots,z_{n'})$ by choosing from the
auxiliary spanning set, with $n'=d_T-n$. Such an extension is
obtained by passing to the quotient $W/\operatorname{span}(y)$
and choosing a basis from the images of the spanning set.
The original family is unchanged. Let
\[
 \mathcal G(\eta)=\det(\langle\Theta_\eta w_i,
                            \Theta_\eta w_j\rangle)_{ij}.
\]
A basis change from an orthonormal basis of $W$ shows
$\mathcal G(\eta)=Z(\eta)\det(\langle w_i,w_j\rangle)$: the Gram
matrix is congruent by its change-of-basis matrix.
In particular $\mathcal G(1)\ne0$ and
$\mathcal G(1)/\mathcal G(\eta)=Z(1)/Z(\eta)$.

For fixed $\iota$ at $\ell$, use matrices $U_\iota$ and
$U_{\bar\iota}$ from the index assumption, applying both at the
\emph{same} twist $\eta=\eta_\iota$, allowed at both
identifications. Form the row transformations on the concatenated
family with blocks
$A_\iota=\operatorname{diag}(U_\iota,\gamma(\eta)I_{n'})$
and $A_{\bar\iota}=\operatorname{diag}(U_{\bar\iota},
\gamma(\eta)I_{n'})$. The auxiliary block need not be a matrix
over $\mathcal O_\iota$: it is its lifted vectors that are integral.
Both transformed lifted families are integral, for their
respective identification. By (\ref{c3:local-pair}) all entries
of their mixed pairing matrix have $\iota$-norm at most one.
Its determinant is
\[
 \overline{\det U_{\bar\iota}}\,\det U_\iota\,
 |\gamma(\eta)|^{2n'}\mathcal G(\eta).
\]
The two index determinants have norms each at least
$\ell^{-e_\ell}$ because conjugation carries $\bar\iota$ back
to $\iota$. Using the ultrametric determinant bound gives
\[
 \big|\,|\gamma(\eta)|^{2n'}\mathcal G(\eta)\,\big|_\iota
       \le\ell^{2e_\ell}.
\]
Hence
\begin{equation}\label{c3:local-height}
 \operatorname{ht}(|\gamma(\eta_\iota)|^{2n'}
                    \mathcal G(\eta_\iota))\le2\Delta .
\end{equation}
The two determinant corrections account for the coefficient $2$ on the index.

At infinity, for each $\sigma$ the Gram determinant of the
concatenated conjugate frame is bounded by the product of the
Gram determinants of its two blocks, by Fischer's inequality
(or by taking the exterior product and projecting the added
vectors to the orthogonal complement of the first span).
Hadamard on the second block gives
\[
 \det\operatorname{Gram}(v_\sigma(w))\le
 e^{c+2n'x_1^2b^2}.
\]
To pass to the conjugate of the scalar $\mathcal G(1)$ we must not
assume that $\sigma$ commutes with conjugation. Put
$\sigma'=\mathrm{conj}\circ\sigma\circ\mathrm{conj}$.
Rationality of $G_{\rm cl}$ gives
$\sigma\mathcal G(1)=\det(v_{\sigma'}(w_i)^*G_{\rm cl} v_\sigma(w_j))$.
The determinant Cauchy--Schwarz inequality on the exterior power,
applied to the two wedge products, bounds its absolute value by
the geometric mean of the corresponding two Gram determinants.
Thus
\begin{equation}\label{c3:arch-final}
 |\sigma\mathcal G(1)|\le e^{c+2n'x_1^2b^2}.
\end{equation}

Factor the same nonzero constant, at each $\iota$, as
\[
 \mathcal G(1)=
 (|\gamma(\eta_\iota)|^{2n'}\mathcal G(\eta_\iota))
 (|\gamma(\eta_\iota)|^{-2})^{n'}
 \frac{Z(1)}{Z(\eta_\iota)} .
\]
Apply the height bounds (\ref{c3:local-height}),
(\ref{c3:gammaheight}), (\ref{c3:chainheight}) and then
(\ref{c3:product}) with (\ref{c3:arch-final}):
\[
 0\le c+2\Delta+n'(2x_1^2b^2+b)+d_T(b+V-H).
\]
There are at most $4d_T/b^2$ auxiliary vectors. Since
$4(2x_1^2+b^{-1})\le b$ by the admissible scales,
\[
 n'(2x_1^2b^2+b)\le d_T\,4(2x_1^2+b^{-1})\le d_T b .
\]
The inequality $b\le V$ then absorbs both terms $d_T b$, proving
\eqref{c3:theorem}.
\end{proof}

The family selected in Theorem B enters this comparison through
its original $n$ ordered columns. Its witnesses provide both the
untwisted norm and the uniform local correction in the final
comparison.

\subsection{Elementary scale checks}\label{c3:scales}

The raw scales give
$L^{100}\le x$, $x^{10^6}\le x_1$, $x_1^{10^6}\le b$,
$\log L\ge10^{12}C_P$ and
$C_P\ge64$. In particular $L>e^{200}$,
$x\ge100$, $b\ge x^{10^{12}}$. The source data give
$\sqrt{x}\le m\le2\sqrt{x}$, $x^5\le r\le3x^5$,
$r/2\le k\le r$, $\delta=5r$, $1\le f\le x^{12}$.
Consequently
\[
 B_0=\frac{f x b L^3}{mk}\ge\frac{f bL^3}{6x^{9/2}}
 \ge100f x^{5000},\qquad
 B_0\ge100x_1.
\]
The last inequality follows directly from
$b/x_1\ge x_1^{10^6-1}\ge x^{10^6(10^6-1)}$.
For the leading local defect term, its ratio to $B_0$ is at most
$4(4\cdot10^6e^{180})/L^2<1/2$; the remaining terms of
the regular-family budget satisfy
$(1+2\delta)x^{999}+x^{4000}+\log G_0+x^{30}\le x^{4001}$
using $\log G_0\le x^{32}$ and $x\ge100$.
The same bound on $\log G_0$ gives the twist-error inequality
$(4f+3)r\log G_0+2\le x^{999}$.
For the archimedean multiplier,
$(fmb/k)/B_0=m^2/(xL^3)\le4/L^3$,
hence $fmb/k+3x^{34}\le B_0$, and
$x^{33}+8f\le B_0$ as well. For the one-lift comparison,
divide its required absorption by $kB_0$:
\[
 \frac{3fmbL^2+rB_0+10mx_1+6kB_0}{kB_0}
 \le\frac{12}{L}+2+\frac{40}{L^3}+6<L .
\]
Also $2m^2\delta x_1\le120x^6x_1\le b$,
$4(2x_1^2+b^{-1})\le12x_1^2\le b$ and
$b\le fxbL^4=V$. Finally
$2x^{201}+x^{99}\le x^{202}$ and
$x^{99}+x^{202}b^2\le2x_1^2b^2$.
These estimates use no dimension lower bound and no sign assumption
on the main amount. The group-order and class-number bounds used in these
inequalities are detailed in Appendix~\ref{sec:finite}.
The auxiliary hook and box estimates enter through Aux.

\subsection{Construction hypotheses for Theorem C}\label{c3:frontier}

Theorem C uses the following assertions for the admissible datum.
They identify where its proof uses a construction beyond the
height and determinant arguments.

\begin{description}
\item[I (target and multiplicity).]
The discrete-series spaces on $W_j$ and $Y_0$ satisfy the
cuspidal trace comparison under the hypotheses of
Appendix~\ref{sec:trace}. That appendix deduces both the common
dimension $d_T$ and equality of the determinants of the common
Hecke operator. The integral pairing, coordinate metric and
triviality of the stabilizers used in the determinant calculation
are derived in \S\ref{c3:section}. Injectivity of the chain uses K.

\item[R (arithmetic at the chosen twists).]
For the rank-$r$ main summands of every relevant eigensystem and
the accompanying rank-one lines, assume the Galois representations,
the simultaneous Selmer and Stickelberger polynomial data, and
the tame realization used in \S\ref{c3:twists}. Assume the
crystalline defect is bounded by the stated envelope and the
first-order congruences yield a clearing operator on each generator
with the indicated valuation, including the rank-one pairing
for all test forms. Appendix~\ref{sec:finite} derives the
root-of-unity specialization, explicit defect-envelope sum, geometric
class bounds and resulting $B_0$ budget from its stated arithmetic
inputs. These conclusions supply items 1--3 of \S\ref{c3:twists}.

\item[G (gluing and exchange).]
For each $j$ the period formula with the nonzero $c_i$ and the
products-spanning assertion in \S\ref{c3:single} hold. Its normalized
scalar quotient is the exchange-value ratio
$\rho_\vartheta=|c'_1|^2/|c'_\vartheta|^2$.
Gluing takes the translated integral frame to integral forms at
every place above $\ell$ coprime to the twist. The glued period
is orthogonal to false lifts, and on that subspace the projection
onto lifts is contractive for every conjugate norm. At trivial
twist the gluing inequality has squared constant $e^{2A_{\rm gl}}$ with
$A_{\rm gl}=fmbL^2$; its class, filling and hook estimates are
$(|B|G_0)^2\le e^{A_{\rm gl}}$ and $C_f,h_j\le e^{x^{999}}$.
The three exchange factors are the critical-value products on the
specified integer intervals, with the indicated character, sign and
main-pair weight conditions. These inputs give
\eqref{c3:compression} and \eqref{c3:compressionheight}, and account
for the $k/2$ selected main points and $5m$ remaining points.

\item[F (rank-one values).]
Assume the common-constant logarithmic norm formula, its flat
integral algebraic parts and conjugate denominator bounds in
\S\ref{c3:single}, and the stated rank-one formula for the
Hecke ratio. Use the valuation-attainment estimate of
Appendix~\ref{sec:jets} for the twelve-dimensional block in
the full degree range. Assume the uniform Taylor scaling
$|\sigma T|\le e^{fmb/k}$ and the four-block identity realizes
the specified pairing on a form integral above $\ell$ and
orthogonal to false lifts. The proof of \eqref{c3:test} details the
Fourier averaging, the prime-power formula and the multiplication
of factors at opposed identifications. Together these assumptions
supply the term $-A_0$ in \eqref{c3:main-height}.

\item[K (Rallis and elementary exchanges).]
For the finite sharp data and archimedean edge types of the
single lifts and of the triangle assume the injectivity and
Petersson coset-sum formula in the range
$2\,\mathrm{rank}(\text{source})<\mathrm{rank}(\text{target})$
where it is used. Assume the algebraic realization of each
Gram pairing, including its index and dimension, the local block
mass formula, and the square commutation and adjoint laws for the
triangle, with the $Q_{ij}$ moduli of \S\ref{c3:chain}. The
parabolic index and twist-factor formulas there fix the residual
norm power independently of $\eta$. The determinant calculation
then gives \eqref{c3:chainfactor} for the entire chain.

\item[Aux (auxiliary path).]
Use the shared short-basis bound of \S\ref{A:source} at the grid
level with the box and conductor bounds in \S\ref{c3:auxiliary}.
Assume the onto map $R$ and normalized first-to-last path-period
identity there, with the product modulus \eqref{c3:gamma};
the shifted delta forms test the indicated dual weight.
The common denominator satisfies $1\le N\le e^{x^{99}b^2}$ and
the conjugate period bound is $e^{x^{200}b^2}\|\sigma v\|\sqrt B$.
Assume the normalized dual Weyl pairing has the integral hook
inverse and isometric archimedean duality stated there, with
$hN e^{x^{200}b^2}\le e^{x^{201}b^2}$; the integral good
projection has the squared-norm cost specified there. The
fixed rescaling, local and archimedean dual recovery, and
completion of the original family then give
\eqref{c3:auxbounds}--\eqref{c3:gammaheight}.
\end{description}

On the constructed family, the local Taylor estimate in F follows from the
curve and frame normalizations in \S\ref{sec:unitrow-frame}--\ref{sec:other-witnesses}.
Section~\ref{rankone:main} treats the rank-one character in R.
Section~\ref{pathmath:dual} derives the dual pairing in Aux, and
\S\ref{pathmath:denom} bounds the common denominator under its one-edge and
period inputs. Each statement is for the normalizations specified there.
